\documentclass[8pt]{extarticle}
\usepackage[margin=0.12in, top=0.10in, bottom=0.10in]{geometry}
\usepackage{amsmath, amssymb}
\usepackage{multicol}
\usepackage{graphicx}
\usepackage{titlesec}
\usepackage{booktabs}
\usepackage{array}
\usepackage{xcolor}
\setlength{\columnsep}{0.26cm}
\setlength{\parindent}{0in}
% --- Squeezing Space ---
\linespread{0.50}
\titlespacing*{\section}{0pt}{0.9pt}{0.3pt}
\titlespacing*{\subsection}{0pt}{0.5pt}{0.2pt}
% -----------------------
\date{}
\begin{document}
\normalsize \textbf{ENGR 232 Final Exam Notesheet} \hfill Sean Balbale \hfill Spring 2026 \vspace{0.03cm}
\hrule
\vspace{0.03cm}
\begin{multicols*}{2}

%==============================================================
\section*{Chapter 2: Atomic Structure \& Bonding}
\subsection*{Key Equations}
\textbf{Average Atomic Weight:} $A_{avg} = \sum f_i A_i$ \quad ($f_i$: fraction of occurrence) \\
\textbf{Force of Attraction:} $F_A = \frac{|Z_1 Z_2| e^2}{4\pi\epsilon_0 r^2}$ \\
($Z_1,Z_2$: ion valences, $r$: interatomic distance, $e=1.602\times10^{-19}$ C, $\epsilon_0=8.85\times10^{-12}$ F/m) \\
\textbf{\% Ionic Character:} $\%IC = \left(1-e^{-0.25(X_A-X_B)^2}\right)\times 100$ \\
($X_A, X_B$: electronegativities)
\subsection*{Concepts}
\textbf{Bohr Model:} Electrons orbit nucleus in discrete energy levels ($E=hf$). \\
\textbf{Bonding:} \textit{Ionic}—electron transfer. \textit{Covalent}—electron sharing. \textit{Metallic}—sea of delocalized electrons.

%==============================================================
\section*{Chapter 3: Crystal Structures}
\subsection*{Key Equations \& Geometry}
\textbf{SC:} $a=2R$, 1 atom/cell. \quad \textbf{BCC:} $a=\frac{4R}{\sqrt{3}}$, 2 atoms/cell. \\
\textbf{FCC:} $a=2R\sqrt{2}$, 4 atoms/cell. \\
\textbf{APF:} $\frac{\text{Vol. of atoms}}{\text{Vol. of unit cell}} = \frac{n(\frac{4}{3}\pi R^3)}{a^3}$ \quad (SC=0.52, BCC=0.68, FCC=0.74) \\
\textbf{Theoretical Density:} $\rho = \frac{nA}{V_c N_A}$ \quad ($N_A=6.022\times10^{23}$ atoms/mol) \\
\textbf{Miller Indices (Planes):} (1) If plane passes through origin, establish new origin. (2) Find intercepts on $x,y,z$. (3) Take reciprocals. (4) Clear fractions. (5) Enclose in $(hkl)$; negatives get a bar.
\subsection*{Concepts}
\textbf{Isotropy/Anisotropy:} Isotropic—properties independent of direction; anisotropic—vary with direction.

%==============================================================
\section*{Chapter 4: Imperfections in Solids}
\subsection*{Key Equations}
\textbf{Total Atomic Sites:} $N = \frac{N_A \rho}{A}$ \\
\textbf{Equilibrium Vacancies:} $N_v = N\exp\!\left(-\frac{Q_v}{kT}\right)$ \quad ($k=8.62\times10^{-5}$ eV/atom·K) \\
\textbf{Mass to Weight \%:} $C_1 = \frac{m_1}{m_1+m_2}\times100$ \\
\textbf{Weight \% to Atom \%:} $C_1' = \frac{C_1/A_1}{C_1/A_1+C_2/A_2}\times100$ \\
\textbf{Atom \% to Weight \%:} $C_1 = \frac{C_1'A_1}{C_1'A_1+C_2'A_2}\times100$
\subsection*{Concepts}
\textbf{Vacancy:} Atom missing from lattice site. \textbf{Self-Interstitial:} Extra atom in void space.

%==============================================================
\section*{Chapter 5: Diffusion}
\subsection*{Key Equations}
\textbf{Fick's First Law (Steady-State):} $J = -D\frac{\Delta C}{\Delta x}$ \\
\textbf{Mass Diffused:} $M = JAt = -DAt\frac{\Delta C}{\Delta x}$ \\
\textbf{Fick's Second Law (Non-Steady-State):} \\
$\frac{C_x - C_0}{C_s - C_0} = 1 - \mathrm{erf}\!\left(\frac{x}{2\sqrt{Dt}}\right)$ \\
\textbf{Diffusion Coefficient:} $D = D_0\exp\!\left(-\frac{Q_d}{RT}\right)$ \quad ($R=8.31$ J/mol·K) \\
\textbf{Activation Energy (Two Temps):} $Q_d = -R\frac{\ln D_1 - \ln D_2}{1/T_1 - 1/T_2}$

\subsection*{erf Table \& Interpolation: $\frac{y-y_1}{y_2-y_1}=\frac{x-x_1}{x_2-x_1}$}
{\scriptsize
\begin{tabular}{@{}ll@{\ }ll@{\ }ll@{\ }ll@{}}
\toprule
$z$ & erf & $z$ & erf & $z$ & erf & $z$ & erf\\
\midrule
0.00&0.000 & 0.30&0.329 & 0.60&0.604 & 1.00&0.843\\
0.10&0.112 & 0.40&0.428 & 0.70&0.678 & 1.20&0.910\\
0.20&0.223 & 0.50&0.521 & 0.80&0.742 & 1.50&0.966\\
0.25&0.276 & 0.55&0.563 & 0.90&0.797 & 2.00&0.995\\
\bottomrule
\end{tabular}}

%==============================================================
\section*{Chapter 6: Mechanical Properties}
\subsection*{Key Equations}
\textbf{Eng. Stress:} $\sigma=F/A_0$ \quad \textbf{Eng. Strain:} $\epsilon=\Delta l/l_0$ \\
\textbf{Shear:} $\tau=F/A_0$, \; $\gamma=\tan\theta$ \\
\textbf{Hooke's Law:} $\sigma=E\epsilon$ \; (elastic only) \quad $\tau=G\gamma$ \\
\textbf{Poisson's Ratio:} $\nu=-\epsilon_x/\epsilon_z = -\epsilon_y/\epsilon_z$ \\
\textbf{True Stress/Strain:} $\sigma_T=\sigma(1+\epsilon)$; \; $\epsilon_T=\ln(1+\epsilon)$ \\
\textbf{Strain Hardening:} $\sigma_T=K\epsilon_T^n$ \\
\textbf{Resilience:} $U_r\approx\frac{1}{2}\sigma_y\epsilon_y=\frac{\sigma_y^2}{2E}$ \\
\textbf{Elastic Elongation:} $\Delta l = \frac{l_0 F}{A_0 E}$
\subsection*{Concepts}
\textbf{Yield Strength:} 0.002 strain offset. \textbf{Tensile Strength:} Max on engineering curve. \textbf{\% Elongation:} $\frac{l_f-l_0}{l_0}\times100$.

%==============================================================
\section*{Chapter 7: Dislocations \& Strengthening}
\subsection*{Key Equations}
\textbf{Resolved Shear Stress:} $\tau_R = \sigma\cos\phi\cos\lambda$ \\
\textbf{Yield Strength:} $\sigma_y = \frac{\tau_{CRSS}}{(\cos\phi\cos\lambda)_{\max}}$ \\
\textbf{Hall-Petch:} $\sigma_y = \sigma_0 + k_y d^{-1/2}$ \\
\textbf{\% Cold Work:} $\%CW = \left(\frac{A_0 - A_d}{A_0}\right)\times100$
\subsection*{Concepts}
\textbf{Slip System:} Slip plane + slip direction. \textbf{Strain Hardening:} Ductile metal gets harder/stronger as plastically deformed (more dislocations).

%==============================================================
\section*{Chapter 8: Failure}
\subsection*{Key Equations}
\textbf{Stress Concentration:} $\sigma_m = 2\sigma_0\sqrt{a/\rho_t} = K_t\sigma_0$ \\
\textbf{Critical Stress (Brittle):} $\sigma_c = \sqrt{\frac{2E\gamma_s}{\pi a}}$ \\
\textbf{Fracture Toughness:} $K_c = Y\sigma_c\sqrt{\pi a}$ \\
\textbf{Cyclic:} $\Delta\sigma=\sigma_{\max}-\sigma_{\min}$; $\sigma_m=(\sigma_{\max}+\sigma_{\min})/2$; $R=\sigma_{\min}/\sigma_{\max}$ \\
\textbf{Creep Rate:} $\dot{\epsilon}_s = K_2\sigma^n\exp(-Q_c/RT)$ \\
\textbf{Larson-Miller:} $m=T(C+\log t_r)$
\subsection*{Concepts}
\textbf{Fatigue:} Structural damage under cyclic stress. \textbf{S-N Curve:} Fatigue lifetime = cycles to failure at given $\sigma$; fatigue strength = $\sigma$ for given cycles.

%==============================================================
\section*{Chapter 9: Phase Diagrams}
\subsection*{Key Equations — Lever Rule}
For a two-phase region with tie line from $C_\alpha$ to $C_\beta$ at composition $C_0$: \\
$W_\alpha = \frac{C_\beta - C_0}{C_\beta - C_\alpha}$ \quad $W_\beta = \frac{C_0 - C_\alpha}{C_\beta - C_\alpha}$ \\
\textbf{Eutectic Microconstituent:} $W_e = \frac{C_0' - 18.3}{43.6}$ (Pb-Sn) \\
\textbf{Fe-C Hypoeutectoid} ($C_0' < 0.76$ wt\%C): \\
$W_p = \frac{C_0' - 0.022}{0.74}$ \quad $W_{\alpha'} = \frac{0.76 - C_0'}{0.74}$ \\
\textbf{Fe-C Hypereutectoid} ($C_0' > 0.76$ wt\%C): \\
$W_p = \frac{6.70 - C_1'}{5.94}$ \quad $W_{Fe_3C'} = \frac{C_1' - 0.76}{5.94}$
\subsection*{Fe-C Key Points}
Eutectoid: $\gamma(0.76\%C) \rightarrow \alpha(0.022\%C) + Fe_3C(6.7\%C)$ at 727°C \\
Eutectic: $L \rightarrow \gamma + Fe_3C$ at 1147°C, 4.30 wt\%C \\
Phases: $\alpha$-ferrite (BCC, soft); austenite $\gamma$ (FCC); cementite $Fe_3C$ (hard) \\
\textbf{Hypoeutectoid} ($<$0.76\%C): proeutectoid $\alpha$ + pearlite \\
\textbf{Hypereutectoid} ($>$0.76\%C): proeutectoid $Fe_3C$ + pearlite
\subsection*{Reactions (Definitions)}
\textbf{Eutectic:} $L \rightarrow \alpha + \beta$ (liquid $\to$ two solids) \\
\textbf{Eutectoid:} $\gamma \rightarrow \alpha + Fe_3C$ (solid $\to$ two solids) \\
\textbf{Peritectic:} $L + \alpha \rightarrow \beta$ (liquid + solid $\to$ solid)
\subsection*{Example: Fe-C Lever Rule (HW7)}
0.45 wt\%C alloy (hypoeutectoid) at room temp: \\
$W_p = \frac{0.45 - 0.022}{0.74} = 0.578$ \quad $W_{\alpha'} = \frac{0.76-0.45}{0.74} = 0.419$

%==============================================================
\section*{Chapter 10: Phase Transformations}
\subsection*{Key Equations}
\textbf{Critical Radius for Nucleation:} \\
$r^* = -\frac{2\gamma T_m}{\Delta H_f}\cdot\frac{1}{T_m - T} = \frac{-2\gamma}{\Delta G_v}$ \\
($\gamma$: surface free energy, $\Delta H_f$: latent heat of fusion, $T_m$: melting temp K, $T$: transformation temp K) \\
\textbf{Activation Free Energy:} \\
$\Delta G^* = \frac{16\pi\gamma^3 T_m^2}{3\Delta H_f^2}\cdot\frac{1}{(T_m-T)^2}$ \\
\textbf{Avrami Equation:} $y = 1 - \exp(-kt^n)$ \\
($y$: fraction transformed, $k,n$: material constants) \\
\textbf{Transformation Rate:} $\mathrm{rate} = 1/t_{0.5}$
\subsection*{Example: Critical Radius (HW9, 10.3)}
Cu: $T_m=1085$°C, nucleates at $849$°C, $\Delta H_f=-1.77\times10^9$ J/m$^3$, $\gamma=0.200$ J/m$^2$: \\
$r^* = \frac{-2(0.200)(1085+273)}{-1.77\times10^9}\cdot\frac{1}{1085-849} = 1.30\times10^{-9}$ m
\subsection*{Example: Avrami Rate (HW9, 10.8)}
$n=2.5$, $y=0.40$ at $t=200$ min. Find $k$: $k = \frac{-\ln(1-y)}{t^n} = \frac{-\ln(0.60)}{200^{2.5}} = 9.0\times10^{-8}$. \\
Find $t_{0.5}$: $t_{0.5} = \left[\frac{-\ln(0.5)}{k}\right]^{1/n} = \left[\frac{0.693}{9.0\times10^{-8}}\right]^{1/2.5} = 311$ min; rate $= 1/311 = 3.2\times10^{-3}$ min$^{-1}$

%==============================================================
\section*{Chapter 12: Ceramics}
\subsection*{Key Equations}
\textbf{Theoretical Density:} $\rho = \frac{n'(\Sigma A_C + \Sigma A_A)}{V_C N_A}$ \\
($n'$: formula units/cell, $A_C,A_A$: cation/anion atomic weights, $V_C$: unit cell volume) \\
\textbf{Frenkel Defect \#:} $N_{fr} = N\exp\!\left(-\frac{Q_{fr}}{2kT}\right)$ \\
\textbf{Schottky Defect \#:} $N_s = N\exp\!\left(-\frac{Q_s}{2kT}\right)$
\subsection*{Concepts}
\textbf{Frenkel Defect:} A cation vacancy + cation interstitial pair (cation moves from site to void). \\
\textbf{Schottky Defect:} A paired set of cation vacancy AND anion vacancy (maintains charge neutrality).
% --- Add to Chapter 12 ---
  \textbf{Ceramic Structures ($n'$):} \\
  Rock Salt (NaCl, FeO, MgO): $n'=4$ \\
  Cesium Chloride (CsCl): $n'=1$ \\
  Zinc Blende (ZnS, SiC): $n'=4$ \\
  Fluorite/Antifluorite: $n'=4$ \\
  Perovskite: $n'=1$ \quad Spinel: $n'=8$

\subsection*{Example: Schottky (HW9, 12.29)}
NaCl at 801°C ($T=1074$ K), $Q_s=2.3$ eV: \\
$N_s/N = \exp\!\left(\frac{-2.3}{2(8.62\times10^{-5})(1074)}\right) = \exp(-12.43) = 3.97\times10^{-6}$

%==============================================================
\section*{Chapter 14: Polymers}
\subsection*{Key Equations}
\textbf{Number-Avg Mol. Weight:} $M_n = \sum x_i M_i$ \quad ($x_i$: number fraction) \\
\textbf{Weight-Avg Mol. Weight:} $M_w = \sum w_i M_i$ \quad ($w_i$: weight fraction) \\
\textbf{Degree of Polymerization:} $DP = M_n / m$ \quad ($m$: repeat unit mol. wt) \\
\textbf{Repeat Unit MW (copolymer):} $m = \sum f_j m_j$ \quad ($f_j$: mole fraction)
\subsection*{Example (HW10)}
Polypropylene $M_n = 10^6$ g/mol, repeat unit C$_3$H$_6$: \\
$m = 3(12.01)+6(1.01) = 42.09$ g/mol; $DP = 10^6/42.09 = 23,760$

%==============================================================
\section*{Chapter 18: Electrical Properties}
\subsection*{Key Equations}
\textbf{Ohm's Law:} $V=IR$ \quad \textbf{Resistivity:} $\rho=RA/l$ \quad \textbf{Conductivity:} $\sigma=1/\rho$ \\
\textbf{Current Density:} $J=\sigma\mathcal{E}$ \quad \textbf{E-field:} $\mathcal{E}=V/l$ \\
\textbf{Metal Conductivity:} $\sigma=n|e|\mu_e$ \\
\textbf{Matthiessen's Rule:} $\rho_{total}=\rho_t+\rho_i+\rho_d$ \\
\textbf{Intrinsic Semiconductor:} $n=p=n_i$; \; $\sigma = n_i|e|(\mu_e+\mu_h)$ \\
\textbf{n-type extrinsic:} $\sigma \cong n|e|\mu_e$ \quad \textbf{p-type:} $\sigma \cong p|e|\mu_h$ \\
\textbf{Capacitance (vacuum):} $C = \epsilon_0 A/l$ \quad $\epsilon_0=8.85\times10^{-12}$ F/m \\
\textbf{Capacitance (dielectric):} $C = \epsilon A/l = \epsilon_r \epsilon_0 A/l$ \\
\textbf{Charge:} $C=Q/V$
\subsection*{Constants}
$|e| = 1.6\times10^{-19}$ C; Si: $\mu_e=0.14$ m$^2$/V·s, $\mu_h=0.048$ m$^2$/V·s, $n_i=1.1\times10^{16}$ m$^{-3}$ \\
Ge: $\mu_e=0.38$, $\mu_h=0.18$ m$^2$/V·s, $n_i=2.4\times10^{19}$ m$^{-3}$
\subsection*{Concepts}
\textbf{Intrinsic Semiconductor:} Conductivity from pure material; $n=p$ (e.g., pure Si, Ge). \\
\textbf{Extrinsic Semiconductor:} Conductivity controlled by dopants. \textit{n-type} = donor impurity (Group V, e.g., P in Si), excess electrons. \textit{p-type} = acceptor impurity (Group III, e.g., B in Si), excess holes.
\subsection*{Example: Intrinsic Carrier Concentration (HW10)}
GaAs: $\sigma=3\times10^{-7}$ ($\Omega\cdot$m)$^{-1}$, $\mu_e=0.80$, $\mu_h=0.04$ m$^2$/V·s: \\
$n_i = \frac{\sigma}{|e|(\mu_e+\mu_h)} = \frac{3\times10^{-7}}{(1.6\times10^{-19})(0.84)} = 2.2\times10^{12}$ m$^{-3}$
\subsection*{Example: Extrinsic Conductivity (HW10)}
Ge doped with $5\times10^{22}$ m$^{-3}$ Sb (n-type), $\mu_e=0.1$, $\mu_h=0.05$ m$^2$/V·s: \\
$\sigma \cong n|e|\mu_e = (5\times10^{22})(1.6\times10^{-19})(0.1) = 800$ ($\Omega\cdot$m)$^{-1}$
\subsection*{Example: Capacitance}
Parallel plate, $A=2\times10^{-4}$ m$^2$, $l=1\times10^{-3}$ m, $\epsilon_r=6.0$: \\
$C = (6.0)(8.85\times10^{-12})(2\times10^{-4})/(1\times10^{-3}) = 1.06\times10^{-12}$ F

%==============================================================
\section*{Chapter 19: Thermal Properties}
\subsection*{Key Equations}
\textbf{Heat Capacity:} $C = dQ/dT$ (J/mol·K) \\
\textbf{Specific Heat:} $c = dQ/(m\,dT)$ (J/kg·K — per unit mass) \\
\textbf{Linear Thermal Expansion:} $\Delta l/l_0 = \alpha_l \Delta T$ \\
\textbf{Volume Expansion:} $\Delta V/V_0 = \alpha_v \Delta T \approx 3\alpha_l\Delta T$ \\
\textbf{Heat Flux (Fourier's Law):} $q = -k\,dT/dx$ \\
\textbf{Thermal Stress:} $\sigma = E\alpha_l(T_0 - T_f)$
\subsection*{Concepts}
\textbf{Heat Capacity ($C$):} Energy required to raise 1 mol of material by 1 K. \\
\textbf{Specific Heat ($c$):} Heat capacity per unit mass; same concept but normalized per kg. Measured at constant pressure ($C_p$) or volume ($C_v$); $C_p > C_v$. \\
\textbf{Thermal Conductivity:} Metals highest (free electrons), ceramics moderate (phonons), polymers lowest. \\
\textbf{Thermal Shock:} Fracture from rapid temp changes; $TSR \propto \sigma_f k/(E\alpha_l)$.

%==============================================================
\section*{HW Problem Reference}
\subsection*{Phase Diagram — Compositions at a Temperature (HW7)}
Given a binary phase diagram at equilibrium, locate the tie line at the temperature. Read compositions of each phase from the phase boundaries, then use the lever rule to get mass fractions.
\subsection*{Resistivity Calculation (HW10)}
$\rho = RA/l$. Example: wire of length $l$, cross-section $A=\pi r^2$, resistance $R$. Find $\rho$ or $\sigma=1/\rho$. \\
Example: $A=\pi(0.001)^2=3.14\times10^{-6}$ m$^2$, $R=0.03$ V/A$=0.003$ $\Omega$ per 300 mm: $\rho=RA/l = (0.003)(3.14\times10^{-6})/(0.3) = 3.14\times10^{-8}$ $\Omega\cdot$m.
\subsection*{Ceramic Density (HW9, 12.16)}
Rock salt (NaCl): $n'=4$ formula units/cell, $A_{Na}=22.99$, $A_{Cl}=35.45$, $a=0.562$ nm: \\
$\rho = \frac{4(22.99+35.45)}{(0.562\times10^{-9})^3(6.022\times10^{23})} = 2.16\times10^3$ kg/m$^3$
\subsection*{Polymer MW (HW10)}
$M_n=\sum x_i M_i$; $M_w=\sum w_i M_i$; $DP=M_n/m$. Use midpoint of each range as $M_i$.
\subsection*{Nucleation (HW9)}
Solve for $k$ from $y=1-\exp(-kt^n)$: rearrange to $k=\frac{-\ln(1-y)}{t^n}$. Then find $t_{0.5}$ from $\ln(0.5)=-kt_{0.5}^n$.
\subsection*{Common Constants}
$R=8.314$ J/mol·K, $k_B=8.62\times10^{-5}$ eV/atom·K, $N_A=6.022\times10^{23}$, $|e|=1.6\times10^{-19}$ C, $\epsilon_0=8.85\times10^{-12}$ F/m

\end{multicols*}
\end{document}